\section{Conclusion}
Experimental results show that Quake reduces query latency compared to baseline approaches under dynamic and skewed workloads, without requiring manual tuning. It achieves high recall, matching the performance of an oracle for setting the query parameter nprobe. Compared to existing partitioned indexes like Faiss and SCANN, Quake reduces query latency by A) adaptively maintaining index partitions and B) maximizing memory bandwidth during query processing. Compared to graph indexes like SVS, HNSW, and DiskANN, Quake offers more efficient indexing and updates while matching or reducing query latency. In summary, our evaluation shows Quake minimizes query latency while meeting recall targets on dynamic workloads with skewed access patterns.

\paragraph{Acknowledgments}
We would like to thank our shepherd, Nitin Agrawal, and the reviewers for their valuable feedback and efforts in making this a stronger paper. This work was supported by NSF grant CNS-2237306, Apple Scholars in AIML PhD Fellowship, and UW-Madison Hilldale Undergraduate Research Fellowship. This work was also supported by DARPA under the grant AIE DARPA-PA-22-01. The U.S. Government is authorized to reproduce and distribute reprints for Governmental purposes notwithstanding any copyright notation thereon. Any opinions, findings, and conclusions or recommendations expressed in this material are those of the authors and do not necessarily reflect the views, policies, or endorsements, either expressed or implied, of DARPA or the U.S. Government.